%==============================================================================
% Synthetic Loan Data Generation and Risk/Approval Modelling
% IEEE-style research paper. Compile with:  pdflatex main.tex  (run twice)
%==============================================================================
\documentclass[conference]{IEEEtran}

\usepackage[utf8]{inputenc}
\usepackage[T1]{fontenc}
\usepackage{graphicx}
\usepackage{booktabs}
\usepackage{amsmath}
\usepackage{array}
\usepackage{url}
\usepackage{xcolor}
\usepackage[hidelinks]{hyperref}

% A compact code/verbatim block for the listings we show.
\usepackage{listings}
\lstset{
  basicstyle=\ttfamily\footnotesize,
  breaklines=true,
  columns=fullflexible,
  frame=single,
  framesep=4pt,
  aboveskip=6pt,
  belowskip=6pt,
}

\begin{document}

\title{Synthetic Loan-Application Data Generation and \\
Machine-Learning Models for Credit-Risk Grading \\
and Approval Decisions}

\author{
\IEEEauthorblockN{Author Name}
\IEEEauthorblockA{Loan Risk \& Approval Project\\
Synthetic Data and Modelling Pipeline\\
\today}
}

\maketitle

%------------------------------------------------------------------------------
\begin{abstract}
Real consumer-lending data is restricted by privacy law and commercial
confidentiality, which makes it difficult to build and openly demonstrate
credit models. This paper presents a fully reproducible pipeline that
\emph{generates realistic synthetic loan-application data} and trains
machine-learning models on it to predict two outcomes: an applicant's
\emph{credit-risk grade} and the lender's \emph{approval decision}. The work
progresses through two iterations: Version~1, a broad United-States--style
dataset with 33 features, and Version~2, a focused dataset grounded in the
Cambodian micro-lending market with 15 features. Both attach
transparent, rule-based labels so that the ground truth is auditable rather
than a black box. On a held-out test set, the risk model attains approximately
91\% accuracy and the approval model approximately 78\%; we show that the gap is
a \emph{deliberate} consequence of injecting realistic stochasticity into the
approval rule, not a modelling deficiency. The deliverable is a config-driven
codebase that cleanly separates data generation from model training and exports
serveable pipelines.
\end{abstract}

\begin{IEEEkeywords}
synthetic data, credit risk, loan approval, machine learning, random forest,
reproducibility, micro-finance.
\end{IEEEkeywords}

%------------------------------------------------------------------------------
\section{Introduction}

Lenders must decide (a)~how risky a borrower is and (b)~whether to approve a
requested loan. Building machine-learning models for these tasks normally
requires historical loan books, which are private. Without data, one cannot
prototype, teach, or interview around the problem. The central question of this
work is therefore: \emph{can we synthesise data that is realistic enough to
support meaningful credit modelling, while keeping the labelling logic fully
transparent?}

\subsection{Objectives}
\begin{enumerate}
  \item Generate a synthetic loan-application dataset that encodes plausible
        socio-economic correlations (income, region, employment, loan
        products).
  \item Attach rule-based labels for a \emph{risk grade} and an \emph{approval
        status}.
  \item Train and evaluate models that recover those labels.
  \item Make the pipeline reproducible and production-shaped, exporting saved
        pipelines that can be served from an application.
\end{enumerate}

\subsection{Contributions}
\begin{itemize}
  \item A documented, seed-reproducible synthetic data generator for the
        Cambodian lending context.
  \item A transparent points-based risk scorecard and a probabilistic approval
        rule that together define the learning targets.
  \item Two trained \texttt{scikit-learn} pipelines and an evaluation harness.
  \item A software design that decouples data generation from training, the two
        communicating only through a documented CSV schema.
\end{itemize}

%------------------------------------------------------------------------------
\section{Background and Motivation}

\subsection{Why synthetic data}
Synthetic data sidesteps privacy and licensing constraints and lets us
\emph{control} the data-generating process. The trade-off is that a model
trained on it only learns the assumptions baked into the generator. We mitigate
this by keeping those assumptions explicit and inspectable in source code rather
than hidden.

\subsection{Why the Cambodian context}
Version~1 used generic United-States--style ranges. Version~2 was re-grounded in
the Cambodian micro-finance market because this makes the synthetic ranges
concrete and defensible: USD-denominated loans, province-level geography,
product types common to Cambodian micro-finance institutions (moto, agriculture,
small-and-medium-enterprise loans), and the high nominal interest rates typical
of the sector (e.g.\ 24--36\% on moto loans). The dataset then reflects a real
domain rather than abstract noise.

%------------------------------------------------------------------------------
\section{The Two Versions}

The project evolved through two iterations, both preserved as reference
notebooks. Table~\ref{tab:versions} summarises the differences.

\begin{table}[t]
\caption{Comparison of the two project versions.}
\label{tab:versions}
\centering
\renewcommand{\arraystretch}{1.25}
\begin{tabular}{p{1.5cm}p{2.6cm}p{2.6cm}}
\toprule
\textbf{Aspect} & \textbf{Version 1 (legacy)} & \textbf{Version 2 (productionised)} \\
\midrule
Market   & US-style, generic & Cambodian micro-lending \\
Features & 33 & 15 \\
Targets  & Approval (class.), Risk score (regression) & Risk category (3-class), Loan status (binary) \\
Risk label & Weighted sub-scores & Points scorecard $\to$ 3 bands \\
Approval & Deterministic threshold & Probabilistic, by risk band \\
Models   & Logistic Reg., RF, Gradient Boosting & Random Forest pipelines \\
Saved?   & No (notebook only) & Yes (\texttt{joblib} + metrics) \\
\bottomrule
\end{tabular}
\end{table}

\subsection{Version 1: broad feature exploration}
Version~1 generated 40{,}000 records with 33 features spanning demographics
(age, education, marital status, dependents), credit behaviour (credit score,
utilisation, inquiries, payment history, bankruptcies), and balance-sheet items
(assets, liabilities, net worth, savings). It derived an interest rate from
credit score, amount and term, computed a total debt-to-income ratio, then
applied a deterministic weighted-sum approval rule and a second weighted
aggregation bucketed into Low/Medium/High risk. It trained Logistic Regression,
Random Forest and Gradient Boosting models, reaching about 91\% approval
accuracy and $R^2 \approx 0.76$ on the risk-score regression. Its weakness was
complexity: 33 features, several weakly motivated, and no model persistence.

\subsection{Version 2: focused, grounded, productionised}
Version~2 reduced the feature set to the 15 most decision-relevant variables and
re-grounded them in Cambodia. It introduced the points-based scorecard for risk
and a probabilistic approval rule, trained one Random Forest pipeline per
target, and saved the pipelines for serving. This is the version rebuilt into
the production package; Version~1 remains a reference.

%------------------------------------------------------------------------------
\section{Methodology}

\subsection{Data-generating process}
Each application is constructed in a fixed \emph{causal order}, so downstream
variables depend on upstream ones, producing realistic correlations rather than
independent noise:
\begin{lstlisting}
Province -> RegionType -> EmploymentType -> Income
                              | Age _____/
                              v
LoanType -> LoanAmount -> Term, Rate -> DTI
                              v
RiskScore -> RiskCategory -> LoanStatus
\end{lstlisting}

The principal conditional rules are: region drives the employment mix (rural
areas skew toward farming, urban toward salaried work); income depends on
region, employment and age (prime age 25--45 receives a higher ceiling); the
loan product depends on the applicant segment; the loan amount is capped at five
times annual income; and term and interest rate are product-specific. The
debt-to-income ratio (DTI) uses the standard amortising-annuity payment
\begin{equation}
A = P \cdot \frac{r(1+r)^{n}}{(1+r)^{n}-1},
\end{equation}
where $P$ is the principal, $r$ the periodic interest rate and $n$ the number of
periods; $\text{DTI} = A / \text{income}$.

\subsection{Labelling}
\textbf{Risk} is a transparent scorecard in $[0,100]$. Each factor contributes
risk points and mitigating conditions subtract them; the maxima are given in
Table~\ref{tab:scorecard}. The score is bucketed into Low ($\le 35$),
Medium ($36$--$65$) and High ($>65$).

\begin{table}[t]
\caption{Risk scorecard factor weights.}
\label{tab:scorecard}
\centering
\renewcommand{\arraystretch}{1.2}
\begin{tabular}{lc}
\toprule
\textbf{Factor} & \textbf{Max points} \\
\midrule
Employment type        & 25 \\
Loan-to-income ratio    & 20 \\
Loan purpose            & 15 \\
Debt-to-income (DTI)    & 15 \\
Credit history          & 12 \\
Region                  & 10 \\
\bottomrule
\end{tabular}
\end{table}

\textbf{Approval} is intentionally \emph{probabilistic}: a base probability set
by risk band (Low 85\%, Medium 45\%, High 15\%) is adjusted by collateral,
credit history, employment and product, after which a Bernoulli draw decides the
outcome. This injects irreducible noise so the task is non-trivial.

\subsection{Modelling}
Both targets are learned by independent Random Forest classifiers wrapped in a
\texttt{scikit-learn} pipeline:
\begin{lstlisting}
ColumnTransformer(
  StandardScaler(numeric) +
  OneHotEncoder(categorical))
-> RandomForestClassifier(
   n_estimators=200, max_depth=12,
   class_weight="balanced")
\end{lstlisting}
Bundling preprocessing inside the pipeline guarantees train/serve consistency:
the saved artifact consumes raw application records directly, with no separate
feature-engineering step that could drift.

\subsection{Experimental setup and leakage control}
We generate 40{,}000 applications and use an 80/20 stratified train/test split
with a fixed seed. We report accuracy, macro and weighted $F_1$, the confusion
matrix, and a per-class report. Crucially, the raw \texttt{RiskScore} (the
intermediate that defines the risk band) and the \texttt{RiskCategory} label are
\emph{excluded} from the feature set; both models consume the same 15 raw
features, so the saved pipelines never observe leakage columns.

%------------------------------------------------------------------------------
\section{Results}

On the held-out test set of 8{,}000 applications, the models achieve the scores
in Table~\ref{tab:results}. A representative risk confusion matrix (rows are
true classes, columns predicted, order Low/Medium/High) is
\[
\begin{bmatrix}
4061 & 278 & 0 \\
171 & 2624 & 192 \\
0 & 56 & 618
\end{bmatrix}.
\]

\begin{table}[t]
\caption{Held-out test-set performance.}
\label{tab:results}
\centering
\renewcommand{\arraystretch}{1.25}
\begin{tabular}{lcc}
\toprule
\textbf{Model} & \textbf{Accuracy} & \textbf{Macro $F_1$} \\
\midrule
Risk (Low/Medium/High)     & $\approx 0.91$ & $\approx 0.89$ \\
Approval (Approved/Reject) & $\approx 0.78$ & $\approx 0.76$ \\
\bottomrule
\end{tabular}
\end{table}

\textbf{Interpretation.} The risk model is highly accurate because its label is
a deterministic function of the features; a tree ensemble recovers a scorecard
well, and its errors concentrate on \emph{adjacent} bands rather than skipping a
band. The approval model plateaus near 0.78 \emph{by design}: the approval label
contains a Bernoulli coin-flip, so part of the outcome is irreducible. A model
scoring near 0.99 here would instead signal label leakage or a broken
deterministic rule.

%------------------------------------------------------------------------------
\section{Discussion}

The project demonstrates a complete, reproducible workflow from data synthesis
to a serveable model, with a disciplined separation between the generation and
training packages, which share only a documented CSV contract. Evaluation is
honest: the approval ceiling is explained rather than hidden. Engineering choices
that support this include a single configuration file (sample size,
distributions, features and hyper-parameters are data, not code),
single-seed reproducibility, and an auditable ground-truth scorecard.

\subsection{Limitations}
The data is synthetic and encodes hand-tuned assumptions; it is not validated
against a real loan book, and the distributions are designed rather than fitted
to survey data. Sensitive-adjacent attributes (gender, province) are present for
realism but no fairness audit has been performed; consequently the models must
not be used for real lending decisions.

\subsection{Future work}
Promising extensions include calibrating distributions against published
Cambodian micro-finance statistics, adding probability-calibrated outputs and
tunable decision thresholds for the approval model, a fairness analysis across
gender and region, gradient-boosting and logistic baselines for comparison, and
unit tests for the scorecard and samplers.

%------------------------------------------------------------------------------
\section{Conclusion}

A carefully designed synthetic dataset can support a credible, end-to-end
credit-modelling workflow without any real customer data. By grounding Version~2
in the Cambodian micro-lending market and keeping the labelling logic
transparent, the pipeline produces models whose strengths (a highly learnable
risk grade) and limits (an intentionally noisy approval task) are both
explainable. The result is a reproducible, production-shaped codebase suitable
for demonstration, teaching, and integration into a web application.

%------------------------------------------------------------------------------
\section*{Reproducibility}
The full pipeline is reproduced with:
\begin{lstlisting}
pip install -r requirements.txt
python -m data_generation.run
python -m model_training.train
python -m model_training.predict
# or simply:  make all
\end{lstlisting}

\begin{thebibliography}{1}
\bibitem{mitchell2019}
M.~Mitchell \emph{et al.}, ``Model cards for model reporting,''
in \emph{Proc. Conf. Fairness, Accountability, and Transparency}, 2019,
pp.~220--229.
\bibitem{sklearn}
F.~Pedregosa \emph{et al.}, ``Scikit-learn: Machine learning in Python,''
\emph{J. Mach. Learn. Res.}, vol.~12, pp.~2825--2830, 2011.
\end{thebibliography}

\end{document}
